\documentclass[11pt]{article}
\usepackage[a4paper,margin=25mm]{geometry}
\usepackage[T1]{fontenc}
\usepackage{lmodern}
\usepackage{booktabs}
\usepackage{amsmath}
\usepackage{graphicx}
\usepackage{hyperref}
\usepackage{xcolor}
\usepackage{listings}
\lstset{basicstyle=\ttfamily\small, columns=fullflexible, keepspaces=true,
  breaklines=true, frame=single, framerule=0.3pt, rulecolor=\color{gray}}
\newcommand{\code}[1]{\texttt{#1}}
\newcommand{\us}{\textmu s}

\title{\code{noisetoy.py}: the expectation for white, independent noise}
\author{BGOTSCAN}
\date{2026-10-15}

\begin{document}
\maketitle

\section{Purpose}

The note on the Antarctica runs (\code{docs/antarctica.tex}) opens with
the structure of the noise and, for each measurement, states first what
white, independent noise would give. This script draws those
expectations from simulated channels with the geometry and the rms of
the data, so that the measured figures of \code{wavefit.py} and
\code{cmnoise.py} can be read against them.

\section{Usage}

\begin{lstlisting}
python3 src/py/noisetoy.py [--sigma 3.8] [--samples 1500] [--channels 64] [--events 236] [--common 0] [--save NAME]
\end{lstlisting}

\begin{center}
\begin{tabular}{l l p{80mm}}
\toprule
option & default & meaning \\
\midrule
\code{--sigma S} & 3.8 & rms per sample, counts (the measured value) \\
\code{--samples N} & 1500 & samples per record (12~\us{} at 8~ns, the tail used in the data) \\
\code{--channels C} & 64 & channels \\
\code{--events E} & 236 & records \\
\code{--common A} & 0 & rms of a component common to all channels, counts, for the third figure \\
\code{--seed} & 1 & random seed \\
\code{--save NAME} & & save the figures as \code{src/py/plots/NAME\_noisetoy\_*.png} \\
\bottomrule
\end{tabular}
\end{center}

\section{Figures}

\paragraph{\code{channel}: one channel.} Gaussian white noise of rms
$\sigma$ per sample. Three panels: the power spectrum, flat up to the
Nyquist frequency of 62.5~MHz, averaged over the records; the
autocorrelation, 1 at zero lag and 0 elsewhere, on the raw samples and
on the 62-sample running mean, where the smoothing spreads the spike
over its width; and the rms of the running mean against its width $N$,
following $\sigma/\sqrt N$, from 1 sample to 5~\us{} (the last points fall below the line because the mean of the
record is removed before smoothing and the running mean is evaluated away from the edges: a finite-record effect,
present identically in the measured figure of \code{wavefit.py}).

\paragraph{\code{channels}: independent channels.} $C$ channels of the
same noise, independent. Two panels: the correlation matrix of their
125-sample running means over the records, ones on the diagonal and
zero elsewhere within the statistical scatter $1/\sqrt{\rm samples}$;
and the rms of the sum of $n$ channels against $n$, following
$\sigma\sqrt n$, for the raw samples and for the running means.

\paragraph{\code{common}: a correlated component.} Only with
\code{--common} above zero: the same two panels with a slow component
of rms $A$ (a 1~\us{} running mean of white noise, scaled) added to
every channel. The matrix shows a uniform off-diagonal value
$A^2/(A^2 + \sigma_{\rm slow}^2)$ and the sum grows faster than
$\sqrt n$, towards $n$. This is what a common noise would look like in
the measured figures.

\section{Examples}

\begin{lstlisting}
python3 src/py/noisetoy.py --save toy            # docs/figs/noisetoy_channel.png, noisetoy_channels.png
python3 src/py/noisetoy.py --common 0.3 --save toyc
\end{lstlisting}

\end{document}
